% Options for packages loaded elsewhere
\PassOptionsToPackage{unicode}{hyperref}
\PassOptionsToPackage{hyphens}{url}
%
\documentclass[
  11pt,
]{article}
\usepackage{amsmath,amssymb}
\usepackage{iftex}
\ifPDFTeX
  \usepackage[T1]{fontenc}
  \usepackage[utf8]{inputenc}
  \usepackage{textcomp} % provide euro and other symbols
\else % if luatex or xetex
  \usepackage{unicode-math} % this also loads fontspec
  \defaultfontfeatures{Scale=MatchLowercase}
  \defaultfontfeatures[\rmfamily]{Ligatures=TeX,Scale=1}
\fi
\usepackage{lmodern}
\ifPDFTeX\else
  % xetex/luatex font selection
\fi
% Use upquote if available, for straight quotes in verbatim environments
\IfFileExists{upquote.sty}{\usepackage{upquote}}{}
\IfFileExists{microtype.sty}{% use microtype if available
  \usepackage[]{microtype}
  \UseMicrotypeSet[protrusion]{basicmath} % disable protrusion for tt fonts
}{}
\makeatletter
\@ifundefined{KOMAClassName}{% if non-KOMA class
  \IfFileExists{parskip.sty}{%
    \usepackage{parskip}
  }{% else
    \setlength{\parindent}{0pt}
    \setlength{\parskip}{6pt plus 2pt minus 1pt}}
}{% if KOMA class
  \KOMAoptions{parskip=half}}
\makeatother
\usepackage{xcolor}
\usepackage[margin=1in]{geometry}
\usepackage{longtable,booktabs,array}
\usepackage{calc} % for calculating minipage widths
% Correct order of tables after \paragraph or \subparagraph
\usepackage{etoolbox}
\makeatletter
\patchcmd\longtable{\par}{\if@noskipsec\mbox{}\fi\par}{}{}
\makeatother
% Allow footnotes in longtable head/foot
\IfFileExists{footnotehyper.sty}{\usepackage{footnotehyper}}{\usepackage{footnote}}
\makesavenoteenv{longtable}
\setlength{\emergencystretch}{3em} % prevent overfull lines
\providecommand{\tightlist}{%
  \setlength{\itemsep}{0pt}\setlength{\parskip}{0pt}}
\setcounter{secnumdepth}{-\maxdimen} % remove section numbering
\ifLuaTeX
  \usepackage{selnolig}  % disable illegal ligatures
\fi
\usepackage{bookmark}
\IfFileExists{xurl.sty}{\usepackage{xurl}}{} % add URL line breaks if available
\urlstyle{same}
\hypersetup{
  pdftitle={Envelope contact on finite graph cylinders, including cycles},
  pdfauthor={Xean research note},
  hidelinks,
  pdfcreator={LaTeX via pandoc}}

\title{Envelope contact on finite graph cylinders, including cycles}
\author{Xean research note}
\date{3 October 2026}

\begin{document}
\maketitle

\section{Envelope contact on finite graph cylinders, including
cycles}\label{envelope-contact-on-finite-graph-cylinders-including-cycles}

Partial theorem, 3 October 2026. Its entire proof received an independent
source-aware review through Xean native Codex, using gpt-6.1-sol with
xhigh reasoning. No earlier proof, review, or computer output is assumed
as a mathematical premise. The exact review and artifact bindings are
preserved with this paper.

\subsection{Result and scope}\label{result-and-scope}

\textbf{Theorem.} Let \(G\subset N_{G,\mathbb{R}}\) be a finite pure
rational polyhedral graph with positive rational edge weights \(w_e\).
Use its actual saturated primitive outgoing directions \(v_{v,e}\) and
primitive edge lengths. Suppose at every vertex \[
 K_v:=\{b\in\mathbb{R}_{\geq0}^{E(v)}:\sum_e b_ev_{v,e}=0\}
       =\mathbb{R}_{\geq0}(w_e)_e.                 \tag{1}
\] Bounded edges, unbounded ends, and arbitrary graph cycles are
allowed. Let \(E\) be a complete rational affine space of dimension
\(m\), with its actual lattice, and give \(X=G\times E\) the direct
product lattice and inherited edge weights. Put \(d=m+1\).

Here real PA means continuous real-affine pieces on a \textbf{finite
rational spatial complex}. For every real-PA PC reference \(\gamma\)
positive at infinity and every bounded real-PA function \(u\), the
actual raw envelope \(P=P_\gamma(u)\) is real PAPC on a finite rational
spatial complex, satisfies \(\gamma-C\leq P\leq\gamma+u\leq\gamma+C\)
for \(C=\|u\|_\infty\), and is decreasing uniformly real-regularizable.
Its full recursive measure has \[
 \operatorname{supp}(P^d\cdot[X])\subset\{P=\gamma+u\},\qquad
 \int_X(P-\gamma-u)\,\operatorname{MA}_{\mathrm{poly},\mathbb{R}}(P)=0,
 \quad \operatorname{MA}_{\mathrm{poly},\mathbb{R}}(P)=P^d\cdot[X]. \tag{2}
\] The functions need not be separable. In particular this includes the
specified triangular graph times every complete affine \(E\), for all
the references and obstacles in the assignment.

This is a special case of original P05.2, not its arbitrary
balanced-pair classification. Higher-dimensional balanced spaces outside
these graph cylinders remain unresolved here. Status remains partial; the metadata flag \texttt{candidate=false}
records that the original classification is not solved. No literal
coefficient conclusion, analytic realization, formal verification,
or priority certification is claimed.

\subsection{Definitions and source
boundary}\label{definitions-and-source-boundary}

The source inspected is Botero--Mazzon--Pille-Schneider,
\emph{Monge--Ampère measures on balanced polyhedral spaces},
arXiv:2603.08664v1, manuscript dated 10 March 2026 (BMPS). Its
Definitions 3.2--3.3 use endpoints \(x_i\in X\cap N_{\mathbb{Q}}\), real
coefficients \(\alpha_i\geq0\), \(\sum_i\alpha_i=1\), and the actual
centre \(x=\sum_i\alpha_i x_i\). A combination is a PCC if on some
finite rational structure there are cells \(\tau\preceq\sigma_i\) with
\(x\in\tau\), \(x_i\in\sigma_i\). A finite raw PC function
\(h:X\to\mathbb{R}\) satisfies every resulting Jensen inequality, on
every permitted structure. Its continuity is not assumed in this proof.

Definition 6.2 gives the raw envelope \[
 P_\gamma(u)(x)=\sup\{h(x):h\in\mathrm{PC}(X),\ h\leq f:=\gamma+u,
                 \ h\leq\gamma+C_h\text{ for some }C_h\in\mathbb{R}\}. \tag{3}
\] The growth class in Definition 3.14 is the two-sided class
\(\gamma+O(1)\). Definition 4.15 means that on every reference affine
cell the slope on every nonzero recession vector is strictly positive.
Definition 3.1 tests PAPC against every nonnegative local balanced cycle
of every positive dimension on every ordinary open subset. We explicitly
extend both affine coefficients and cycle weights to real numbers.
Definition 3.23's uniform PAPC regularization and Definition 4.4's full
recursive zero-cycle are used in this real convention. Our envelopes
have the explicit regularization \(P+1/j\), whose recursive measures are
identical to that of \(P\).

The literal convention in Definition 2.28, integral tangent slopes and
rational constants, is different. On a fixed finite rational structure
the real coefficients satisfying continuity form a rational linear
system, so our real PA class is also the real span of that structure's
literal PA functions. We do not infer literal regularity.

The source's Proposition 6.13 invokes BBGS22 Proposition 5.13 for PC
locality. Its proof, the controlling book's already established
one-dimensional graph theorem (including cycles), the earlier tree
draft, and the finite-dual proposal are comparison material only. None
is a dependency of the theorem below. The source's automatic raw-PC
continuity, PAPC/PC equivalence proof, and larger intersection-theory
dependency chain are not imported. All convex, graph, rational-domain,
and intersection facts needed here are supplied below. The controlling
catalogue and appendix were inspected to delimit the result; no P06.2
claim is made.

\subsection{Elementary polyhedral and convex
facts}\label{elementary-polyhedral-and-convex-facts}

We record proofs of the finite-dimensional facts used subsequently. A
cone defined by finitely many rational inequalities has finitely many
rational generators. Remove its rational lineality first. In the pointed
quotient the sum of the defining nonnegative forms is positive away from
zero; its section of value one is bounded, since an unbounded sequence
normalized on the sphere would give a nonzero vector of value zero. A
bounded polyhedron is the convex hull of its finitely many vertices:
within the minimal face of any nonvertex point, a line through the point
ends on smaller-dimensional faces, and induction reduces to vertices.
Vertices solve full-rank subsystems, hence are rational for rational
data. Lift back and add both signs of a rational lineality basis.

Homogenization applies this assertion to a pointed rational polyhedron
and gives a nonempty finite point set and a finite ray set with \[
 \sigma=\operatorname{conv}(A_\sigma)+\operatorname{cone}(R_\sigma),
                    \qquad A_\sigma,R_\sigma\text{ rational}.       \tag{4}
\] The same assertions, without rationality of coordinates, hold for
real inequality data. Eliminating one variable from nonstrict finite
inequalities compares every lower bound with every upper bound,
retaining conditions independent of the variable; a one-sided family
imposes no further condition. Induction proves that projections of
polyhedra are polyhedra. Coefficients of variables remain rational if
they start rational, even if the constant terms are real. In particular
finitely generated convex-hull-plus-cone sets are closed. These
arguments prove the finite-polyhedron and elimination facts used here
rather than invoking an external polyhedral theorem.

All separations used below follow from nearest points in
finite-dimensional Euclidean space. A point outside a nonempty closed
convex set has a nearest point by restricting to a sufficiently large
compact ball. Differentiating squared distance along a segment in the
set gives a separating normal. For a disjoint compact convex set and
closed convex set, distance is positive and attained by the same compact
restriction; differentiation in both sets strictly separates them by
their nearest-point difference. For a cone the finite supremum of the
separating functional is zero. This proves the point/cone and
compact-polytope/cone separations used below, without importing a
separation theorem.

A finite rationally convex function on \(\mathbb{Q}^n\) has a unique
continuous convex completion to \(\mathbb{R}^n\). Corner values bound it
above on rational boxes. Reflection about a rational centre gives a
lower bound on a smaller box: \(F(y)\geq2F(c)-F(2c-y)\). Choose bounds
\(L\leq F\leq M\) on a rational box \(c+[-2R,2R]^n\). For rational
\(y,z\in c+[-R/2,R/2]^n\), set \(\delta=\|z-y\|_\infty\). If
\(\delta>0\), the rational factor \(T=R/\delta\geq1\) puts \(y+T(z-y)\)
in the larger box. Jensen gives \(F(z)-F(y)\leq(M-L)\delta/R\);
reversing the pair gives the opposite bound. This proves local Lipschitz
completion; density proves convexity and uniqueness. A finite convex
function on all of \(\mathbb{R}^n\) is continuous by the same box
argument.

It has an affine support at each point. Separate a point just below its
closed epigraph by the normal at a nearest epigraph point; nearest
points exist by compact restriction, and the separating inequality
follows by differentiating squared distance along line segments. The
vertical normal coefficient is positive, since a nonzero horizontal
normal cannot bound all of \(\mathbb{R}^n\). Normalize it. Evaluating the
resulting affine minorants at \(x\pm Re_i\) bounds their slopes as the
below-epigraph point approaches \((x,F(x))\). A convergent subsequence
yields a support touching at \(x\). This also covers \(n=0\).

For a compact polyhedron \(D\), its outward normal cone at \(p\) is the
cone of the active defining inequality normals. Indeed membership in
this cone implies the normal inequality. If a normal were outside that
closed finitely generated cone, separation would provide a direction
satisfying the active inequalities and increasing the proposed normal; a
sufficiently small step satisfies all inactive inequalities too, a
contradiction. This proves both the assertion and its rational-generator
version when inequality normals are rational, irrespective of real
constants.

Finally, if \(a(p)=\min_i(c_i+g_i(p))\) is concave polyhedral, its
supergradient set at \(p\) is the convex hull of the active \(g_i\). The
directional derivative is their minimum. Separation proves the
supergradient assertion: an alleged supergradient outside that convex
hull gives a violating direction. For a finite sum of such functions
with nonnegative coefficients, the supergradient set is the
corresponding Minkowski sum, directly by expanding the finite minima.
These facts hold also at zero coefficients.

\subsection{Primitive-normal algebra, including actual
indices}\label{primitive-normal-algebra-including-actual-indices}

For every rational cell \(\sigma\), use the saturated tangent lattice
\(N_\sigma=N\cap\operatorname{span}_{\mathbb{R}}(\sigma-\sigma)\). For a
facet \(\tau\prec\sigma\), an inward normal \(n_{\sigma/\tau}\) is a
lift of the primitive positive generator of \(N_\sigma/N_\tau\). For a
real PA \(a\) and a balanced real \(k\)-cycle \(z\), the divisor
coefficient is \[
 (a\cdot z)(\tau)=\sum_{\sigma\succ\tau}z(\sigma)a_{\sigma,\mathrm{lin}}(n_{\sigma/\tau})
 -a_{\tau,\mathrm{lin}}\left(\sum_{\sigma\succ\tau}z(\sigma)n_{\sigma/\tau}\right).       \tag{5}
\] The second sum is tangential by balancing. Changes of normal lift
cancel because the affine pieces agree on \(\tau\). The formula is
bilinear; ambient affine functions have zero divisor. It uses only
incident germs and is therefore local on ordinary opens.

Here are the needed refinement, balancing, and commutativity proofs. At
a new tested \((k-1)\)-face of a refinement, consider its minimal old
face. If that face has dimension \(k\), the two inherited \(k\)-pieces
have the same weight and linear part and opposite primitive quotient
normals, so cancel. If its dimension is \(k-1\), saturated tangents and
incident primitive normal classes agree with the old ones, and (5)
agrees. If its dimension exceeds \(k\), all inherited \(k\)-weights
vanish there. These exhaust the cases. Inherited balancing follows from
the same cancellation at new facets. Thus refinement preserves divisors
and balancing.

To prove balancing of the new divisor and commutativity, fix a
\((k-2)\)-face, quotient its saturated tangent lattice, and subtract
ambient extensions of the common tangential linear parts of the
functions. The local cycle becomes a two-dimensional fan, possibly in a
larger ambient space. Rational hyperplane cuts make its two-cones
pointed. For a cone bounded by primitive rays \(r,s\), let \[
 \chi_{rs}=[N_{C_{rs}}:\mathbb{Z}r+\mathbb{Z}s],\qquad \omega_{rs}=z(C_{rs}).
\] Extending \(r\) to a basis of the actual saturated plane lattice
shows that the primitive inward normal modulo \(\mathbb{R}r\) is
\(s/\chi_{rs}\). No index is replaced by one. Quotient tangent lattices
are saturated: a multiple belonging to the rational tangent subspace
forces the original class to belong there too. Balancing at a ray says
\[
 \sum_{C_{rs}\succ r}(\omega_{rs}/\chi_{rs})s=t_r r.
\] The first divisor has ray coefficient \[
 D_a(r)=\sum_{C_{rs}\succ r}(\omega_{rs}/\chi_{rs})a(s)-t_r a(r).
\] Substituting the balancing equation and pairing the unordered
two-cones proves \(\sum_rD_a(r)r=0\): the terms \(a(s)r+a(r)s\) cancel
exactly. The second divisor at the quotient origin equals \[
 \sum_{\{r,s\}}(\omega_{rs}/\chi_{rs})[a(s)b(r)+a(r)b(s)]
                  -\sum_rt_ra(r)b(r),                       \tag{6}
\] which is symmetric. This proves balancing and commutativity in every
germ, for real coefficients and weights. Formula (5), refinement, and
these identities construct the entire recursive cycle, not just its
maximal-face interior.

Direct product lattices split primitive normals. Hence a function pulled
back from one factor has divisor equal to its divisor in that factor
times the other cycle. As a normalization check, for
\(\sum_{i=1}^d|s_i|\) in lattice-coordinate affine \(d\)-space of weight
\(c\), expansion and commutativity give \(2^d d!c\) at the origin: each
coordinate occurs once. We keep the full recursive normalization, with
no division by \(d!\).

\subsection{Scalar graph facts proved without PC
locality}\label{scalar-graph-facts-proved-without-pc-locality}

If the presentation has a full affine-line edge without vertices, first
subdivide it at a rational point into two oriented ends. This bivalent
subdivision preserves the carrier, weights, actual lattices and PCC tests
by the subdivision argument below. Thus every end strip used subsequently
has a finite base.

The strictly positive vector in (1) implies that the full kernel of the
direction map is one-dimensional. Otherwise a small perturbation of that
vector by another kernel vector stays positive and contradicts (1). At a
bivalent vertex the actual primitive directions are opposite and the
weights equal. At any vertex of larger valence there is no opposite
pair: such a pair would give a second nonnegative relation. Subdivision
at rational edge points therefore preserves (1).

A convex polyhedron contained in \(G\times E\) projects to a convex
subset of \(G\). A straight segment cannot pass a branching vertex,
since there is no opposite pair there. Thus the graph projection is in a
single straight edge chain through bivalent vertices, or is a vertex.
This applies to every permitted rational structure, not just our chosen
subdivision.

\textbf{Scalar lemma.} Suppose \(k\) is finite on rational graph points
and satisfies rational graph PCC inequalities with rational
coefficients. It has a continuous completion, is convex on every edge,
has finite outgoing derivatives \(s_{v,e}\), and satisfies
\(\sum_e w_es_{v,e}\geq0\). Conversely a continuous function convex on
the edges with these inequalities, when added to any affine function of
\(E\), satisfies every actual cylinder PCC, with real coefficients and
all real centres.

For the first assertion, complete the rational Jensen function on each
edge interior. A one-sided limit at an endpoint is finite and at most
the specified endpoint value: an endpoint secant bounds it above and an
interior secant bounds it below. For rational small \(t>0\), the points
\(v+t v_{v,e}\) with coefficients \(w_e/\sum w_e\) are an actual vertex
PCC. Its inequality and the endpoint upper bounds force all limits to
equal \(k(v)\). Every outgoing derivative is bounded above by a finite
secant. Dividing the same PCC inequality by \(t\) prevents any
derivative being \(-\infty\), and gives the weighted derivative
inequality.

For the converse, edgewise convexity and the bivalent inequalities give
ordinary convexity on every straight chain. At a branching centre \(v\),
group the graph displacements of any cylinder PCC as \(t_i v_{v,e}\),
with \(t_i\geq0\), using the preceding structural observation. Their
weighted sum is zero, so the group coefficients
\(b_e=\sum_{i\text{ in }e}\alpha_it_i\) equal \(cw_e\), \(c\geq0\).
Convexity gives \(k(q_i)-k(v)\geq t_i s_{v,e}\). Summing gives a
nonnegative result. At a centre in a straight chain the assertion is
ordinary convexity. Terms with zero graph displacement and the affine
\(E\) terms are handled by the barycentre identity. This proves every
actual PCC directly, including remote endpoints, irrational centres, and
structures crossing bivalent subdivision points. No open-cover locality
theorem is used.

For a raw cylinder competitor \(h\), its rational slice at a rational
graph point \(q\) has the finite convex completion just proved. For a
real dual covector \(p\in E^*\), put \[
 k_{h,p}(q)=\inf_{z\in E_{\mathbb{Q}}}(h(q,z)-p(z))\in[-\infty,\infty).             \tag{7}
\] It satisfies rational graph PCC inequalities: choose rational \(z_i\)
approaching the endpoint infima, apply the actual cylinder PCC at
\((q,\sum_i\alpha_i z_i)\), subtract \(p\), and take the infimum.
Rational coefficients ensure that this last affine coordinate is
rational. If one endpoint has infimum \(-\infty\), vary that endpoint
alone.

If (7) is finite at one rational graph point, it is finite at every
rational point in that connected component. From an edge-interior finite
centre, any other rational point on that edge can be an endpoint in a
rational combination by choosing the other endpoint on the opposite
side. At a finite vertex, any rational remote point \(v+t_ev_{v,e}\) on
a chosen incident edge belongs to a vertex PCC with positive rational
coefficients proportional to \(w_e/t_e\), choosing positive rational
distances within all other incident edges. No endpoint can have value
\(-\infty\). Propagate across the finite connected graph. This works
through cycles and all unbounded ends. The scalar lemma then applies to
(7).

\subsection{Finite partial-infimum data, with tails and recession
constraints}\label{finite-partial-infimum-data-with-tails-and-recession-constraints}

Work in one connected component, identify \(E\) rationally with its
lattice vector space, and take a common finite rational affine structure
for \(f=\gamma+u\) and \(\gamma\), refined by the original edge strips.
Each cell is pointed: a nonzero lineality direction would give two
opposite reference recession slopes, incompatible with positivity at
infinity. Formula (4) therefore applies.

For each cell ray with zero graph coordinate, write its affine obstacle
slope as \(s_r\) and its \(E\)-coordinate as \(r_E\), and define \[
 D=\{p\in E^*:p(r_E)\leq s_r\text{ for every such ray}\}.                    \tag{8}
\] All inequality normals are rational; constants may be real. On every
recession ray \(u\) has zero slope, by boundedness, so (8) is the same
for \(f\) and \(\gamma\). It contains 0 in its interior, because every
nonzero vertical ray has strictly positive reference slope. For \(m=0\),
\(D=\{0\}\) in the zero-dimensional dual space.

Set \(F_p(q)=\inf_{z\in E}(f(q,z)-p(z))\). For \(p\in D\) this is finite
at every \(q\) and its infimum is attained at a vertex of some cell
fibre. Indeed, the fibre is pointed and its recession cone is generated
by the zero-graph cell rays: graph recession is zero on bounded strips
and nonnegative on unbounded strips, so a zero-graph combination cannot
use any positive-graph ray. The objective is nonnegative on this
recession cone, and (4) for the fibre gives a minimizing vertex. There
are finitely many cells. If \(p\notin D\), a violating rational ray
based at a rational point gives \(F_p(q)=-\infty\) at a rational graph
point. Thus (8) is precisely the common finite domain required here.

The set \(D\) is compact. At any fixed rational graph vertex, the
finite-PA function \(f(q,z)\) is bounded above by \(M\|z\|+B\), for some
finite \(M,B\) (use finitely many affine pieces and the dual norm). If
\(\|p\|>M\), a direction with \(p(z)>M\|z\|\) makes the infimum
\(-\infty\). Finiteness for \(p\in D\) bounds \(D\). This proof does not
require \(f\) convex.

Subdivide each bounded edge at the graph coordinates of all its spatial
cell vertices, and each end at those coordinates and its original base.
This is finite and rational. Between consecutive coordinates a fibre
vertex lies on a spatial one-face: the minimal face containing it loses
one dimension under the slicing hyperplane, so has dimension at most
one. A one-face traversing the interval has rational form
\(z(t)=z_i+t b_i\), and obstacle height real-affine in \(t\). Its
feasibility across that open interval is constant, because changes occur
only at projected spatial vertices. Therefore on each such interval \[
 F_p(t)=\min_i\{A_i(p)+t B_i(p)\},                      \tag{9}
\] where the finite index set is independent of \(p\); \(A_i,B_i\) are
affine in \(p\), with rational linear coefficients and real constants.
At a node \(v\), slicing the cells directly gives \(F_p(v)\) as a finite
minimum of affine functions of \(p\), again with rational gradients and
real constants.

For clarity, (9) extends concavely and continuously to each finite
interval endpoint. Its candidate fibre vertices have finite limits
there, so their limiting objective values are at least \(F_p(v)\).
Conversely, fixing any finite \(z\) and using continuity of \(f(t,z)\)
gives \(\limsup_{t\to v}F_p(t)\leq f(v,z)-p(z)\); then take its infimum
over \(z\). This proves equality of the limits. It also proves
branchwise continuity at a graph vertex. The possibility of a different
endpoint fibre type is not discarded.

After the last projected vertex on an end, translate its primitive
parameter to start at 0. Formula (9) holds for all \(t\geq0\), and \[
 \theta_e(p):=\lim_{t\to\infty}F_p(t)/t=\min_i B_i(p),\qquad
 F_p(t)-\theta_e(p)t\text{ is nondecreasing}.             \tag{10}
\] Every \(B_i-\theta_e\geq0\) proves the monotonicity term by term. In
particular the line from any value at most \(F_p(0)\) with slope
\(\theta_e(p)\) lies below the entire tail obstacle. All recession
constraints and finite endpoints have now been retained.

\subsection{The graph incidence polyhedron and its greatest
assignment}\label{the-graph-incidence-polyhedron-and-its-greatest-assignment}

Let \(V\) be the finite node set just obtained. For fixed \(p\in D\),
introduce node values \(a_v\), and impose \[
 a_v\leq F_p(v),\qquad
 L_v(a,p):=\sum_{e=vw\text{ bounded}}\frac{w_e}{\ell_e}(a_w-a_v)
                +\sum_{e\text{ tail at }v}w_e\theta_e(p)\geq0.             \tag{11}
\] Here \(\ell_e>0\) is the actual rational primitive length. All
bounded edges are present in this one system, including each edge of
every cycle. Do not integrate a graph direction into a larger lift or
unwind the graph.

There is a feasible assignment. The raw PC function \(\gamma-C\),
\(C=\|u\|_\infty\), lies below \(f\). Its partial infimum is finite for
\(p\in D\), by (8) applied to \(\gamma\), and the scalar lemma applies.
Replace its restriction on each bounded interval by its endpoint chord,
and on each tail by the line of slope \(\theta_e(p)\) from its endpoint
value. Convexity and the upper bound \(F_p\) show that the chords lie
below the concave data (9); the tail bound follows from (10). The chord
outgoing slopes at both endpoints are at least the old outgoing
derivatives. On a tail the asymptotic upper bound by (9) forces every
secant slope of the old convex function to be at most \(\theta_e(p)\),
so the new outgoing slope is also at least the old one. Thus all
inequalities in (11) hold.

The componentwise maximum of two feasible assignments is feasible. At a
node attaining its value from one assignment, each neighboring value has
only increased, so \(L_v\) has not decreased; the caps also persist.
Choose a feasible witness \(b\). Maximizing \(\sum_v a_v\) after taking
componentwise maximum with \(b\) restricts to the compact box
\(b_v\leq a_v\leq F_p(v)\), and hence attains a maximum. A feasible
assignment exceeding that maximum in any coordinate would increase this
objective by taking componentwise maximum. Consequently there is a
unique componentwise greatest assignment, denoted \(a_v(p)\).

It also dominates every scalar function \(k\) obtained from (7) that
lies below \(F_p\). The chord-and-tail replacement just used dominates
\(k\) everywhere by convexity and the tail slope bound, and its node
vector is feasible for (11). Therefore it is at most the interpolation
of the greatest assignment. Define that interpolation as \(k_p(q)\):
endpoint-linear on every bounded interval and \(a_v(p)+t\theta_e(p)\) on
every tail. It is continuous, scalar PC by the scalar lemma, and
\(k_p\leq F_p\) everywhere.

Every \(a_v(p)\) is a finite concave piecewise-affine function on \(D\),
with rational gradients and real constants. To see this without an
unproved parametric-LP assertion, expand all caps \(F_p(v)\) into their
finite affine minima. Expand each tail minimum \(\theta_e\) in the last
inequalities of (11): the weighted sum of minima is the minimum over all
choices of tail pieces. Thus (11) is one finite linear system in
\((a,p)\) with rational variable coefficients and real constants,
together with (8). Eliminate all \(a_w\), \(w\ne v\), by the elementary
elimination procedure. For fixed \(p\in D\), its section in \(a_v\) is
exactly \((-\infty,a_v(p)]\), because simultaneously lowering all node
values preserves the Laplacians and caps. Projected inequalities
involving \(a_v\) must therefore be upper bounds; zero-coefficient
inequalities concern \(D\). There is at least one upper bound from the
original cap. Dividing positive rational coefficients gives \[
 a_v(p)=\min_j(c_{v,j}+g_{v,j}(p)),\qquad g_{v,j}\text{ rational}.             \tag{12}
\] This also proves continuity on all of \(D\), including its boundary.
Tail slopes already have the same property. No assertion that the dual
vertices are rational is made.

\subsection{Actual raw envelope and finite rational spatial
domains}\label{actual-raw-envelope-and-finite-rational-spatial-domains}

Define \[
 Q(q,z)=\sup_{p\in D}\{k_p(q)+p(z)\}.                                    \tag{13}
\] It is finite: \(D\) is compact and all node values and tail slopes
are continuous. For each fixed \(p\), \(k_p(q)+p(z)\) satisfies every
actual PCC by the scalar lemma, so its supremum does too
\(take the supremum of each Jensen inequality\). Since \(k_p\leq F_p\),
\(Q\leq f\).

For every raw competitor \(h\leq f\), and every rational \((q,z)\), take
a touching affine support of the convex completion of the rational slice
\(h(q,\cdot)\). Its covector \(p\) makes (7) finite and
\(h(q,z)=k_{h,p}(q)+p(z)\). Propagation proves finiteness throughout the
component; a violating ray would contradict \(h\leq f\), so \(p\in D\).
The preceding greatest-assignment argument gives \(k_{h,p}\leq k_p\).
Hence \(h(q,z)\leq Q(q,z)\) at all rational carrier points.

The function \(Q\) is continuous. Partition \(D\) into the common finite
polyhedral chambers of all affine minima in (12) and (10). Each chamber
is compact. On it \(k_p(q)+p(z)\) is affine in \(p\) on a fixed graph
interval, so its maximum is attained at a chamber vertex. Let \(J\) be
the finite union of all these vertices, with no rationality assumption
on their coordinates. Formula (13) on each edge strip is the maximum of
the corresponding finitely many real-affine functions of \((q,z)\). At
graph nodes the formulas agree. This proves continuity on all of \(X\).

To extend the raw upper bound to a real point \(x\), choose any finite
rational structure and its minimal cell containing \(x\). In a
positive-dimensional cell, shrinking simplices with rational vertices in
its relative interior enclose \(x\); they can be obtained by rationally
perturbing the vertices of a small enclosing simplex in the rational
affine hull. Their real barycentric coefficients give actual same-cell
PCCs. Raw Jensen and continuity of \(Q\) imply \(h(x)\leq Q(x)\). A
zero-dimensional cell point is already rational. No raw-competitor
continuity has been used.

Apply this to \(h=\gamma-C\). Together with \(Q\leq f\leq\gamma+C\), it
gives the two-sided growth bound. Consequently \(Q\) itself belongs to
the competitor set (3), and all its competitors are at most \(Q\). Thus
\[
                        P_\gamma(u)=Q.                                  \tag{14}
\] Both the global lower competitor and global upper comparison are
exact, and no local raising or unproved growth gluing is invoked.

Real-affine maxima alone would not prove rational spatial domains. Here
there is an additional rational-gradient mechanism. Fix \(p_0\in J\).
Let \(A_v(p_0)\) be the supergradient polytope of \(a_v\) there, and let
\(B_e(p_0)\) be that of \(\theta_e\). Their vertices are rational by
(12). Let \(N_D(p_0)\) be the outward normal cone of \(D\), generated by
rational normals by (8).

For any finite concave polyhedral \(b\), the condition that \(p_0\)
maximize \(b(p)+p(z)\) over \(D\) is \[
                    z\in-\partial^+b(p_0)+N_D(p_0).                      \tag{15}
\] For sufficiency, a supergradient \(g\) with \(g+z\in N_D\) gives the
global maximum inequality. For necessity, if \((z+\partial^+b)\cap N_D\)
is empty, separate the compact polytope from the closed polyhedral cone.
The separating direction \(r\) has \((g+z)(r)>0\) for every
supergradient and has \(n(r)\leq0\) for every normal. It is a feasible
tangent direction to \(D\); a sufficiently small step is feasible. The
directional derivative, the minimum of these positive quantities,
increases the objective, a contradiction. This supplies the exact
optimality argument, including boundary faces.

On a bounded interval \(t\in[0,\ell]\), \(b=(1-t/\ell)a_v+(t/\ell)a_w\),
so (15) describes the full maximization domain of the affine plane
indexed by \(p_0\) as \[
 z\in-(1-t/\ell)A_v(p_0)-(t/\ell)A_w(p_0)+N_D(p_0).                       \tag{16}
\] It is the convex hull of the rational endpoint point sets
\((0,-\operatorname{Vert}A_v)\), \((\ell,-\operatorname{Vert}A_w)\),
plus the cone of \((0,n)\) for rational generators of \(N_D\). On a tail
its domain is \[
 t\geq0,\qquad z\in-A_v(p_0)-tB_e(p_0)+N_D(p_0),                         \tag{17}
\] generated by the rational points \((0,-\operatorname{Vert}A_v)\) and
rational rays \((1,-\operatorname{Vert}B_e)\), \((0,n)\). These
descriptions remain exact at \(t=0,\ell\), at zero coefficients, and at
degenerate chambers. They prove that every maximization domain is a
rational spatial polyhedron, though its affine value has real
coefficients.

The maximization domains of a finite affine maximum, including all their
faces and removing duplicates, form a complex: intersections are
equality faces of the respective maxima regions. Its faces are rational
because these regions are rational polyhedra. At a graph node (16)--(17)
have exactly the same fibre domain for each \(p_0\), independent of the
neighboring interval, so the strip complexes glue. This gives a finite
rational spatial complex for \(P\) on the actual carrier. Moreover (16)
has vertices only at its two graph endpoints, and (17) only at its base;
faces introduce no new vertices. Thus every possible zero-dimensional
domain face of this complex lies over a node of \(V\). Edge interiors
cannot carry a recursive zero-cycle coefficient. This conclusion
includes all bounded-edge interiors and all tails, rather than omitting
their lower-face terms.

\subsection{PAPC against all local real cycles and actual
regularization}\label{papc-against-all-local-real-cycles-and-actual-regularization}

Let \(z\geq0\) be any local real balanced cycle of any positive
dimension on an ordinary open \(U\subset X\), and refine to make \(P\)
affine. At a tested facet \(\tau\), choose a rational
\(x\in U\cap\operatorname{relint}\tau\). The cone of possible
nonnegative incident weights balanced modulo \(N_\tau\) is a rational
polyhedral cone, so is generated by rational vectors, by the elementary
cone proof above. The divisor coefficient is linear in those weights. It
suffices to test a nonnegative rational incident vector \(c_\sigma\).

Write \(S=\sum c_\sigma>0\),
\(B=\sum c_\sigma n_{\sigma/\tau}\in N_{\tau,\mathbb{Q}}\), and use the
rational inward vectors \(n_{\sigma/\tau}-B/S\). They need not be
integral normals; they have the same positive quotient directions. For
sufficiently small rational \(\varepsilon>0\), the rational points
\(x_\sigma=x+\varepsilon(n_{\sigma/\tau}-B/S)\) belong to
\(\sigma\cap U\). With coefficients \(c_\sigma/S\) their barycentre is
\(x\), giving an actual PCC. Its Jensen inequality is \[
 0\leq\sum c_\sigma[P(x_\sigma)-P(x)]
       =\varepsilon\left(\sum c_\sigma P_\sigma(n_{\sigma/\tau})-P_\tau(B)\right).
\] This is exactly (5). The zero-weight case is zero.
Rational-generation then proves the inequality for every nonnegative
real incident vector, and hence for every real local cycle. Thus \(P\)
is PAPC in the full required sense. Balancing proved above permits
recursion, and every successive divisor is nonnegative. Its full top
cycle is a finite nonnegative atomic measure on the rational domain
vertices. The sequence \(P+1/j\) is PAPC in the same two-sided growth
class and decreases uniformly to \(P\); its divisors agree with those of
\(P\). This establishes the actual operator domain and regularization
without invoking source regularity assertions.

\subsection{Two complete recursive vanishing
lemmas}\label{two-complete-recursive-vanishing-lemmas}

\textbf{Quotient vanishing.} On a local \(d\)-cycle, suppose that after
subtracting an ambient affine function the PA germ is pulled back by a
rational linear map to a function on a complete rational fan in a space
of dimension \(q<d\). Its full top divisor is zero. Let \(T_j\) be the
union of quotient fan faces of codimension at least \(j\). The first
divisor vanishes off the preimage of \(T_1\), since the function is
ambient affine there. Inductively, near the relative interior of a
quotient face of codimension \(j-1\), the preceding cycle is supported
in that face's preimage, where the function is affine on the entire
supporting subspace. Extend this affine function to the ambient space;
balancing makes its next divisor zero. Thus the \(j\)-th divisor is
supported in the preimage of \(T_j\). Since \(T_{q+1}\) is empty, the
top divisor vanishes. Rational common refinements and the local formula
(5) justify this argument on every lower face, without transversality or
a top-face-only measure formula.

\textbf{Lineality vanishing.} If the local carrier fan and residual
function-domain fan have common nonzero rational lineality, choose their
common generalized fan retaining that lineality. Every face contains it;
there is no zero-dimensional face. Every recursive divisor lives on that
fan, so the top zero-cycle is zero. Refinement cannot change this by the
proved refinement identity. Ambient affine terms can be removed first.

\subsection{Full contact at cyclic junction
fibres}\label{full-contact-at-cyclic-junction-fibres}

Let \(x=(v,z_0)\) be any possible atom of the finite rational spatial
complex for \(P\). By (16)--(17), \(v\in V\) and \(z_0\) is rational.
Suppose \(P(x)<f(x)\). Define its compact exposed dual face \[
 C_x=\{p\in D:a_v(p)+p(z_0)=P(v,z_0)\}.                                  \tag{18}
\] It is nonempty and convex. Indeed \(a_v\) is concave, and the
displayed quantity never exceeds its maximum. Formula (12) describes
(18) by inequalities with rational normals \(g_{v,j}+z_0\) and real
constants, together with (8). Its affine-hull direction is therefore
rational: the affine hull is cut out by the inequalities that are
equalities throughout the polyhedron, as seen by choosing a
relative-interior point and varying in its remaining feasible
directions. Their homogeneous coefficients are rational. On \(C_x\),
\(a_v(p)=P(x)-p(z_0)\) is affine.

For each incident bounded edge \(e=vw\), directed away from \(v\), and
each tail, put \[
 s_e(p)=\frac{a_w(p)-a_v(p)}{\ell_e}\quad\text{or}\quad s_e(p)=\theta_e(p).
                                                                         \tag{19}
\] On \(C_x\) these are concave continuous polyhedral functions: \(a_v\)
is affine there. The vertex inequality says \[
                           \sum_e w_es_e(p)\geq0.                         \tag{20}
\] The local germ of \(P-P(x)\) in the outgoing strip with coordinates
\(t\geq0, \zeta=z-z_0\) is exactly \[
                  h_e(t,\zeta)=\max_{p\in C_x}\{p(\zeta)+t s_e(p)\}.       \tag{21}
\] To justify this, use the fixed finite set \(J\) of common dual
chamber vertices. All inactive planes at \(x\) have strictly positive
gap, and therefore stay inactive in a sufficiently small neighbourhood.
Active planes have their dual points in \(C_x\). Conversely every point
of \(C_x\) belongs to a chamber and is a convex combination of its
vertices; equality in the exposed maximum forces those vertices to be
active, and on that chamber all terms in (19) are affine. Thus (21)
follows exactly. This is a local finite-germ identity, not an assumption
of raw continuity or joint-convex reconstruction.

If \(\dim C_x<m\), let \(T\subset E\) annihilate the rational direction
of \(C_x\). It is nonzero and rational. Subtract the ambient affine
\(p_0(z)\), with \(p_0\in C_x\). Formula (21) is invariant under \(T\)
in every strip; the carrier star has \(E\), hence \(T\), as lineality. A
rational function-domain fan retaining \(T\) exists: each coarsest
linearity region is a convex union of finitely many rational tangent
cones of the already constructed spatial complex, and hence is a
rational cone; its equations are invariant under \(T\). Quotient these
cones by \(T\) and lift their common rational refinement. The lineality
lemma proves that the full coefficient at \(x\) is zero. This case
includes all lower-dimensional exposed faces.

Now suppose \(\dim C_x=m\). For \(p\in\operatorname{int}C_x\), choose
\(\eta>0\) with \(p+\eta\overline B^*\subset C_x\). Formula (18) and
(13) give, for every \(y\in E\), \[
 P(v,y)-p(y)\geq P(v,z_0)-p(z_0)+\eta\|y-z_0\|.                           \tag{22}
\] Since \(f(v,\cdot)\geq P(v,\cdot)\) is continuous, \(f(v,y)-p(y)\) is
coercive and attains its infimum. If \(a_v(p)=F_p(v)\), that infimum
equals the first term on the right of (22). The strict inequality away
from \(z_0\) then forces the minimum at \(z_0\) and
\(f(v,z_0)=P(v,z_0)\), a contradiction. Therefore \[
                         a_v(p)<F_p(v)\quad(p\in\operatorname{int}C_x).    \tag{23}
\]

The greatest-assignment property now implies \(L_v(a(p),p)=0\). If this
Laplacian were positive, raise \textbf{only its scalar LP coordinate}
\(a_v\) by a sufficiently small \(\delta>0\). The cap persists by (23).
At \(v\) the Laplacian decreases by \(\delta\sum_{e=vw}w_e/\ell_e\), and
at every neighboring node it increases; tail slopes are fixed. Thus
feasibility persists, contradicting coordinatewise maximality. If there
are no bounded incident edges, raising the coordinate does not change
any Laplacian and gives the contradiction even without a positive
Laplacian; in that case (23) itself is impossible. This is a finite
scalar feasibility argument; it is not a local perturbation of \(P\)
requiring global PC gluing.

Continuity extends the equality to all of \(C_x\): \[
                              \sum_e w_es_e=0\quad\text{on }C_x.           \tag{24}
\] Since each \(s_e\) is concave and all weights are strictly positive,
every \(s_e\) is affine on \(C_x\). Indeed the nonnegative Jensen
deficits of the individual concave functions have weighted sum zero and
must each vanish. On the full-dimensional \(C_x\), choose an interior
point off the finitely many affine-piece walls. Formula (12) and the
rational lengths show that these affine expressions have rational
\(p\)-gradients. Write \[
                    s_e(p)=\alpha_e+p(\beta_e),\qquad\beta_e\in E_{\mathbb{Q}}.
                                                                         \tag{25}
\] The constants \(\alpha_e\) may be real. Equation (24) says
\(\sum w_e\alpha_e=0\) and \(\sum w_e\beta_e=0\).

Let \(S_v\) be the actual rational span of the primitive star
directions. Its only relation is (1). Consequently there is a real
linear form \(A:S_v\to\mathbb{R}\) with \(A(v_{v,e})=\alpha_e\), and a
rational linear map \(B:S_v\to E\) with \(B(v_{v,e})=\beta_e\).
Rationality follows by solving the rational direction matrix for the
rational values, not by assuming ray-generated lattice saturation.
Formula (21) is therefore the restriction to the entire star of \[
                       A(s)+\operatorname{supp}_{C_x}(\zeta+B(s)).        \tag{26}
\] After removing \(A+p_0\circ(\zeta+B(s))\), it factors through the
rational map \((s,\zeta)\mapsto\zeta+B(s)\) to the \(m\)-dimensional
complete space \(E\). The support function of \(C_x-p_0\) has a complete
rational normal fan, because (18) has rational inequality normals; its
normal cones are generated by active normals as proved above. Apply
quotient vanishing with \(q=m<d=m+1\). The entire recursive coefficient
at \(x\) is zero, with actual primitive indices and every intermediate
lower-face contribution retained.

The proof covers every possible atom: original cyclic junctions,
inserted bivalent nodes, all fibre domain faces, and tail bases.
Non-node edge interiors have already been excluded by the
rational-domain generator argument. No tree branch decomposition, cycle
cutting, globally affine lift, or independent support extension is used.
For \(m=0\) the same argument reduces to the known graph contact
mechanism; the new assertion is every \(m>0\) with arbitrary
nonseparable data and cycles. Since the measure is finite nonnegative
and every noncontact coefficient is zero, (2) follows. The integrand is
bounded by the two-sided growth inequalities, and no boundary atom is
introduced: the PA operator is the finite zero-cycle on \(X\), and its
explicit constant-shift regularization has exactly the same measure.
Infinity and all ends have been handled in (8)--(11).

For finitely many connected components, construct \(D\), the LP, and
\(P\) separately. A PCC and a local cycle stay in the component of their
incident cell; the finitely many measures add and a uniform growth
constant is available. This finishes the stated graph-cylinder theorem.

\subsection{Exact triangle and comparison to residual
scope}\label{exact-triangle-and-comparison-to-residual-scope}

For \(A=(0,0)\), \(B=(1,0)\), \(C=(0,1)\), the oriented bounded
primitive directions are \(a=(1,0)\), \(b=(-1,1)\), \(c=(0,-1)\), all of
primitive length one. The leaves \(c-a=(-1,-1)\), \(a-b=(2,-1)\),
\(b-c=(-1,2)\) are primitive. At each vertex the three outgoing vectors
sum to zero and have rank two. Their only relation is the strictly
positive unit vector. Their local lattice spans are saturated: each star
contains a pair of determinant \(\pm1\). Thus this graph meets (1) with
the actual lattices, and the theorem applies to \(G\times E\) in every
dimension.

The tree amalgamation failure is still genuine as a diagnostic: before
imposing the cycle period, \(a,b,c\) are independent and the formal walk
ends at \(a+b+c\ne0\); closing the actual graph imposes \(a+b+c=0\). The
incidence system (11) retains all three bounded edges and never uses
that invalid lift. Its local contact proof also never groups remote
generators into independent branches.

The controlling book already includes cycles in dimension one. Repeating
that theorem is not the contribution. The result here is the full
finite-graph-cylinder subclass in every affine dimension, all
positive-growth real-PA PC references and bounded real-PA obstacles on
finite rational spatial domains, including nonseparable functions. It
addresses the concrete implementation assignment, while original P05.2's
arbitrary balanced-space characterization remains open. The tree draft,
prior two-vertex proof, and finite-dual proposal are not assumed
correct. No additional smoothness hypothesis, separability hypothesis,
regularity-failure alternative, or isolated negative-atom diagnostic
substitutes for the exact contact theorem.

\subsection{Actual nonseparable cyclic example and exhaustive finite
recursion}\label{actual-nonseparable-cyclic-example-and-exhaustive-finite-recursion}

On the specified triangle times \(\mathbb{R}_z\), choose \[
 \gamma(x,y,z)=\frac{|x|+|y|}{10}+|z|.
\] This is the restriction of an ambient convex function and hence
satisfies every actual PCC. Its end slopes in the graph direction are
\(1/5,3/10,3/10\), and its vertical slopes are positive after the
\(|z|\) subdivision; thus it is positive at infinity, including mixed
recession directions. On each of the three A-incident edges let \(t\) be
primitive distance from \(A\), and put \[
 u(t,z)=-\frac32\max\{0,1-|z-1|-2t\};
\] put \(u=0\) on the other edges. On the bounded A-incident edges
\(t=1\) at their opposite endpoints, where the formula vanishes, so the
function is globally continuous. All walls and endpoints are rational,
and \(-3/2\leq u\leq0\). It is nonseparable: its values at
\((t,z)=(0,1),(1/2,1),(0,0),(1/2,0)\) have nonzero rectangular
difference.

Insert \(D,H,J\) at distance \(1/2\) from \(A\) along \(AB,AC\), and A's
leaf, respectively. The dual domain is \([-1,1]\). For
\(0\leq t\leq1/2\) on an A-incident edge with graph reference slope
\(s\in\{1/10,1/10,1/5\}\), fibre vertices occur at \(z=0,2t,1,2-2t\).
Comparing their values of \(f-pz\), for \(-1\leq p\leq1\), gives \[
 F_p(t)=\min\{st,(s+3)t-\tfrac12-p\}.
\] The candidates at \(2t\) and \(2-2t\) are no smaller than \(st\). For
\(t\geq1/2\), \(u=0\), so the infimum is the graph reference value.
Consequently the node caps are \[
 F_p(A)=\min(0,-\tfrac12-p),\quad F_p(B)=F_p(C)=\tfrac1{10},\quad
 F_p(D)=F_p(H)=\tfrac1{20},\quad F_p(J)=\tfrac1{10}.
\] The greatest feasible values of the complete six-node cyclic system
are \[
\begin{split}
 a_A&=\min(0,-\tfrac12-p),\\
 a_B=a_C&=\min(\tfrac1{10},-\tfrac15-p),\\
 a_D=a_H&=\min(\tfrac1{20},-\tfrac15-\tfrac p2,-\tfrac7{20}-p),\\
 a_J&=\min(\tfrac1{10},-\tfrac25-p).
\end{split}                                                       \tag{27}
\] Here is a direct greatest-assignment check, not just an LP execution.
For every feasible vector, bivalent inequalities give
\(a_D\leq(a_A+a_B)/2\), \(a_H\leq(a_A+a_C)/2\), and
\(a_J\leq a_A+1/10\). The B and C inequalities and those bounds imply
\(2a_B-a_C\leq a_A+3/10\), \(2a_C-a_B\leq a_A+3/10\). Substitution
yields \(a_B,a_C\leq a_A+3/10\); combine with their caps and A's cap,
then bound D,H,J. These are exactly the coordinate upper bounds in (27).
The displayed vector attains all those bounds simultaneously. Directly,
its Laplacian slacks at A,B,C are respectively \((2/5,1/5,1/5)\) for
\(p\leq-1/2\), \((7/5+2p,-3/10-p,-3/10-p)\) for \(-1/2\leq p\leq-3/10\),
and \((4/5,0,0)\) for \(p\geq-3/10\); all other slacks are zero. Caps
and feasibility follow. This proves maximality for every \(p\),
including the entire cycle.

Use the exact formula (13), interpolating (27) on all six bounded
intervals and the fixed slopes on the three tails. The common dual
vertices are \[
                         -1,\ -\tfrac12,\ -\tfrac3{10},\ 1.
\] This is an explicit finite formula for the actual envelope by the
theorem proved above. It is itself nonseparable: \(P(D,0)-P(A,0)=1/20\),
whereas \(P(D,1)-P(A,1)=3/20\).

\texttt{cyclic-example.py} evaluates the full local two-dimensional
recursion in exact rational arithmetic. For each node fibre vertex it
first uses (21). Sort the active dual points, insert the branch rays of
slope \[
 \zeta/t=-\frac{s_e(p_{i+1})-s_e(p_i)}{p_{i+1}-p_i},
\] and the shared vertical primitive rays. Multiply rational branch
coordinates to get actual primitive integer rays. Every neighboring ray
pair forms a two-cone in the actual saturated branch plane. Compute its
actual determinant index \(\chi_{rs}\). For each shared ray, compute the
tangential multiple \(t_r\) from \(\sum(\omega_{rs}/\chi_{rs})s=t_rr\),
and then \[
 D_h(r)=\sum_{C_{rs}\succ r}(\omega_{rs}/\chi_{rs})h(s)-t_rh(r),\qquad
                         (h^2\cdot[X])\{0\}=\sum_r D_h(r)h(r).            \tag{28}
\] These are exactly the already proved divisor formulas (5)--(6). The
program verifies balancing and nonnegativity of every first-divisor ray,
rather than replacing the recursion by a fibre-volume formula. All
indices happen to be one for this particular example; the algorithm
retains them. The complete finite output is:

\begin{longtable}[]{@{}lrrrr@{}}
\toprule\noalign{}
Fibre point & \(P\) & \(f\) & \(P-f\) & Full recursive mass \\
\midrule\noalign{}
\endhead
\bottomrule\noalign{}
\endlastfoot
\(A,0\) & 0 & 0 & 0 & \(2/5\) \\
\(A,1\) & \(-1/2\) & \(-1/2\) & 0 & \(58/25\) \\
\(B,0\) and \(C,0\), each & \(1/10\) & \(1/10\) & 0 & \(6/25\) \\
\(B,1\) and \(C,1\), each & \(-1/5\) & \(11/10\) & \(-13/10\) & 0 \\
\(D,0\) and \(H,0\), each & \(1/20\) & \(1/20\) & 0 & 0 \\
\(D,1/2\) and \(H,1/2\), each & \(-1/5\) & \(11/20\) & \(-3/4\) & 0 \\
\(D,1\) and \(H,1\), each & \(-7/20\) & \(21/20\) & \(-7/5\) & 0 \\
\(J,0\) & \(1/10\) & \(1/10\) & 0 & 0 \\
\(J,1\) & \(-2/5\) & \(11/10\) & \(-3/2\) & 0 \\
\end{longtable}

These are all fourteen possible coefficients: the node supergradient
slopes in (27) give exactly these fibre vertices, and (16)--(17) exclude
additional graph-interior vertices. Thus the finite computation is
exhaustive for the explicitly described candidate, with mathematical
envelope identification and spatial coverage supplied separately above.
The four positive atoms have total \(2/5+58/25+2(6/25)=16/5\), and the
complete integral is exactly zero. Independently, the reference's graph
degree is \(1/5+3/10+3/10=4/5\); direct product expansion gives
reference top mass \(2!(4/5)\cdot2=16/5\). Neither mass agreement nor
execution alone establishes contact; the table checks every noncontact
coefficient and the theorem proves the general assertion.

\subsection{Reproducibility and verification
limits}\label{reproducibility-and-verification-limits}

\texttt{diagnostics.py} and \texttt{diagnostic-input.json} provide exact
rational checks of the triangle stars and a three-node cyclic incidence
toy system; \texttt{diagnostic-output.json} records the output. The toy
system is explicitly not asserted to be the partial infimum of an actual
obstacle. It checks incidence arithmetic only.
\texttt{cyclic-example.py} and \texttt{cyclic-example-output.json}
provide the actual nonseparable example and full recursion just
displayed. \texttt{P05.2-partial.tex} is a standalone rendering of this
argument; \texttt{P05.2-partial.pdf} was compiled twice.
\texttt{build-evidence.json}, \texttt{source-hashes.json},
\texttt{environment.txt}, \texttt{commands.txt},
\texttt{author-self-audit.md}, and \texttt{research-notes.md} record
environment, source provenance, build checks, and limitations. Rerun
\texttt{python3\ diagnostics.py}, \texttt{python3\ cyclic-example.py},
and \texttt{python3\ build.py}; the last command renders the TeX and
compiles the PDF twice.

No supervised heavy-compute facility was available in the exposed tools;
none was needed or launched. All computation is small exact arithmetic
or document compilation. Successful execution, source hashes, and
compilation are not mathematical verification. The saved publication
source was opened in the native desktop editor and its built-in compiler
returned success; two local pdfLaTeX passes exported the same source.
The complete independent exact-scope PASS review is included in the
review receipt. That review is not formal verification or a solution of
the original classification.

Bibliography: BMPS, arXiv:2603.08664v1, definitions 2.18--2.37,
3.1--3.3, 3.14, 3.23, 4.4, 4.15, 6.2 and the proof of Proposition 6.13,
inspected for conventions only. The authoritative controlling chapter
and appendix, the known local rank-one graph proof, the tree draft, and
the finite-dual proposal are source/scope comparisons only. No external
mathematical theorem is essential to the proof above.

\end{document}
